% --- SECCIÓN 1 ---
\subsection{Desarrollo Móvil Multiplataforma}

\subsubsection{Framework Flutter}
Flutter es un framework gratuito y de código abierto que creó Google para hacer apps que se compilan de forma nativa y funcionan en varias plataformas, todo desde un solo código base \cite{FlutterDocES}. Lo principal de su idea es cambiar cómo se desarrolla software, ya que te deja armar, probar y lanzar apps para celulares, la web, computadoras de escritorio y hasta dispositivos integrados, usando la misma fuente. Esto ahorra un montón de tiempo y recursos. Básicamente, atiende a lo que se necesita hoy: llegar a los usuarios en cualquier pantalla sin el lío y el gasto de tener códigos diferentes para cada sistema operativo.

Flutter básicamente se arma con dos partes clave: un SDK, que es como un kit de herramientas para desarrollar software, y un framework que usa widgets. El SDK trae lo necesario para convertir el código en lenguaje nativo, ya sea para procesadores ARM o Intel, o incluso en JavaScript, lo que asegura que las apps corran rápido y se sientan suaves para el usuario \cite{FlutterStudyDocker}. El framework está hecho en Dart y viene con una gran colección de widgets que puedes ajustar a tu gusto, dándote control total sobre cada detalle de la interfaz. Así, es más fácil crear diseños que se adapten bien y se vean consistentes en cualquier dispositivo o plataforma. Además, una cosa genial es el \textit{"hot reload"}, que deja ver los cambios en el código casi de inmediato sin que la app pierda su estado actual, y eso acelera mucho el proceso de probar y arreglar cosas.

\subsubsection{Lenguaje Dart}
Dart es un lenguaje de programación actual y con tipos estáticos que Google inventó para crear apps que corran bien en varias plataformas. Su estructura combina una forma de escribir código clara y conocida, parecida a la de Java o C\#, con funciones más avanzadas que lo hacen perfecto para armar interfaces de usuario modernas. Como dice, “Dart es un lenguaje optimizado para aplicaciones UI, enfocado en el desarrollo en cualquier plataforma” \cite{DartWiki2024}. Y por eso, se convierte en la base principal de tecnología para el framework Flutter.

La estructura técnica de Dart permite una compilación que se adapta bien, y eso es esencial para desarrollar en varias plataformas. Puedes compilar el código Dart de antemano (AOT) a código nativo para apps móviles y de escritorio que necesitan correr rápido y eficiente, o hacerlo en el momento (JIT) para cosas como el \textit{``hot reload''} mientras estás programando \cite{DartEvolution2024}. Otra cosa importante es su sistema de tipos con \textit{sound null safety}, que evita los errores típicos de referencias nulas. Lo explican como “un sistema de tipos que garantiza que ninguna expresión cuyo tipo no admita el valor nulo pueda tenerlo” \cite{DartDocHiberus}, y esto hace que el código sea más sólido en general.

\subsubsection{Programación Asíncrona}

La programación asíncrona es algo básico en el desarrollo de apps móviles actuales con Flutter, porque deja correr tareas que podrían tardar mucho, como pedir datos por internet o acceder a bases de datos, sin trabar el hilo principal donde corre todo. Y esto es clave para que la interfaz de usuario se mantenga rápida y sin trabas, ya que "las operaciones asíncronas permiten que su programa complete el trabajo mientras espera que finalice otra operación" \cite{DartAsyncAwait}. Como Dart funciona con un solo hilo, este sistema se basa en un bucle de eventos (\textit{event loop}) y en aislamientos (\textit{isolates}) para manejar tareas al mismo tiempo, donde los hilos se comunican pasando mensajes en vez de compartir memoria directamente \cite{RedwerkAsyncFlutter}.

En Dart, la asincronía se maneja básicamente con objetos Future y las palabras clave async y await. Un \textit{Future<T>} es como una promesa de un valor de tipo T que va a estar listo en algún punto más adelante, y puede estar en proceso o ya terminado, ya sea con el valor correcto o con un error \cite{DartAsyncLanguage}. Gracias a async y await, puedes escribir código que lidia con estas operaciones de manera simple y clara, casi como si fuera código normal que corre en secuencia. Si marcas una función con \textit{async}, esta devuelve automáticamente un Future, y dentro de ella usas \textit{await} para detener temporalmente la ejecución hasta que el Future termine, lo que deja libre el hilo principal para hacer otras cosas mientras espera \cite{DartAsyncAwait}.

% --- SECCIÓN 2 ---
\subsection{Arquitectura de Software}

\subsubsection{Domain-Driven Design (DDD)}

Domain-Driven Design, o DDD, es un método para diseñar software que Eric Evans presentó, y se enfoca en manejar la complejidad al poner el énfasis en el corazón del dominio y su lógica. Básicamente, modela el programa para que represente de manera precisa el mundo del negocio, basado en lo que saben los expertos en esa área \cite{Evans2003DDD}. Lo principal que busca es reducir la distancia entre cómo funciona el negocio en la realidad y el código que se escribe. Para lograrlo, fomenta que los técnicos y los expertos del dominio trabajen juntos de forma creativa, y así pulan un modelo conceptual que todos compartan \cite{WikipediaDDD}. Dos elementos clave en el lado estratégico son el Lenguaje Ubicuo, que es como un vocabulario compartido que usa todo el equipo de manera consistente y que se ve reflejado directamente en el código, y los Contextos Delimitados, o Bounded Contexts, que marcan límites claros donde un modelo del dominio aplica y se mantiene coherente, para evitar confusiones \cite{NetmentorDDD}.

En el aspecto más práctico, DDD ofrece patrones para poner en marcha el núcleo del dominio. Por ejemplo, las Entidades son objetos que tienen una identidad única y que duran con el tiempo, los Objetos-Valor son inmutables y se definen por sus características, y los Agregados son grupos de entidades y objetos-valor que se tratan como una sola unidad para mantener la consistencia, con una raíz que controla el acceso \cite{Evans2003DDD}. Además, incluye otros patrones como Servicios de Dominio, Repositorios y Eventos de Dominio, que ayudan a armar un modelo completo y que exprese bien las ideas \cite{NetmentorDDD}. Este enfoque pone el dominio por encima de los detalles técnicos, y resulta especialmente útil en sistemas complicados con reglas de negocio enredadas, porque genera software que se ajusta mejor a lo que el negocio necesita, es más fácil de mantener y está listo para cambiar con el tiempo \cite{WikipediaDDD}.

\subsubsection{Gestión de Estado Reactiva}
La gestión de estado reactiva es un patrón clave en el diseño de interfaces de usuario actuales, donde la pantalla se refresca sola cada vez que cambia algo en el estado de la app. Lo principal aquí es que la vista actúa como una función que responde al estado, y sincroniza de forma eficiente la interfaz con los datos que hay debajo. Por ejemplo, usando la idea de una bombilla inteligente, "si la variable de datos (la luz ambiental) cambia, la aplicación reacciona y se actualiza automáticamente" \cite{MediumReactivity2024}. Este estilo resuelve cuestiones antiguas como tener que manejar manualmente el DOM o lidiar con desajustes entre los datos y la interfaz, y así mejora la experiencia para el usuario y hace que los desarrolladores sean más productivos \cite{DevToReactiveState}.

Dentro del mundo de Flutter, este patrón se pone en práctica con varias bibliotecas y métodos. Opciones como Provider, Riverpod y Bloc ofrecen formas de tener una sola fuente de verdad para el estado, y solo notifican y reconstruyen los widgets que dependen de una parte específica de los datos cuando esta se modifica \cite{TrioLibraries2026}. Elegir una biblioteca depende de cosas como cuán complicada sea la app, si necesita crecer bien o si prefieres un código más declarativo. En general, este enfoque es vital para armar apps en Flutter que sean fáciles de mantener y que rindan alto, porque ordena el flujo de datos de un modo predecible y simplifica encontrar y arreglar errores \cite{MediumReactivity2024}.

\subsubsection{Patrón Repositorio}

El patrón Repositorio es una capa que se coloca entre la lógica de dominio de la app y las fuentes de datos, para que no se mezclen. Lo principal es separar bien la lógica de la aplicación de cómo funcionan las bases de datos o las APIs, ofreciendo una interfaz siempre igual y limpia para acceder a los datos, y dejando el núcleo de la app sin enterarse de detalles técnicos como si es una API REST o una base local \cite{ModyoRepoPattern}. En Flutter, este patrón va en la capa de datos y se encarga de aislar los modelos del dominio de los orígenes reales, convirtiendo los DTOs en entidades validadas y bien tipadas que la capa de dominio puede usar sin problemas \cite{MediumCleanArchRepo}.

Para implementarlo, primero defines un contrato abstracto con una interfaz o clase abstracta donde listas las operaciones disponibles, y luego creas clases concretas que meten la lógica real, como llamadas HTTP o consultas a bases de datos. Esta abstracción da mucha flexibilidad, ya que puedes cambiar o combinar fuentes de datos (por ejemplo, una API remota o una base local) sin tocar la lógica de negocio ni la interfaz de usuario, y además facilita las pruebas unitarias usando repositorios simulados o fakes \cite{CodewithAndreaRepoPattern, MediumCleanArchRepo}.

% --- SECCIÓN 3 ---
\subsection{Geolocalización y Sistemas GIS}

\subsubsection{Geolocalización GPS}
La geolocalización por GPS es esencial en apps móviles para obtener coordenadas del dispositivo. En Flutter se implementa principalmente con el plugin geolocator, que abstrae las APIs nativas (FusedLocationProviderClient en Android y CLLocationManager en iOS) para obtener la última posición conocida, la actual o actualizaciones continuas \cite{PubDevGeolocator}.

Antes de usar los datos hay que solicitar permisos con Geolocator.requestPermission(), y luego llamar métodos como getCurrentPosition() \cite{ReaniSolucionesGPS}. El plugin también calcula distancias entre puntos geográficos, útil para funcionalidades de proximidad \cite{PubDevGeolocator}. Para un funcionamiento correcto se deben declarar permisos específicos en AndroidManifest.xml y en Info.plist de iOS, manejando todo de forma responsable para respetar la privacidad y garantizar fiabilidad \cite{ReaniSolucionesGPS}.

\subsubsection{Google Maps SDK}
El plugin google\_maps\_flutter es el SDK oficial de Google para meter mapas interactivos de Google Maps directamente en apps Flutter, tanto en Android, iOS como web. Se encarga solo de todo lo pesado: conectar con los servidores de Google, mostrar el mapa y responder a los toques o gestos del usuario, así que reduces un montón de trabajo para tener mapas buenos \cite{GoogleDocsMapsConfig}.

Para usarlo solo agregas la dependencia google\_maps\_flutter al pubspec.yaml y pones una clave API de Google Maps que configures para cada plataforma. Lo mejor es que puedes poner encima del mapa lo que quieras: marcadores para señalar puntos, líneas o polígonos para dibujar rutas o zonas, y mover la cámara del mapa como te dé la gana con código \cite{GoogleCodelabsMaps}. Gracias a que Google Maps cubre casi todo el planeta con actualizaciones constantes, sirve perfecto para cosas simples como mostrar una ubicación o para sistemas más avanzados como seguimiento de repartidores o rutas en tiempo real dentro de una app Flutter \cite{GoogleDocsMapsConfig}.

% \subsubsection{Bases de Datos en Tiempo Real}
% Una base de datos en tiempo real está pensada para que cualquier cambio que se haga se refleje al instante en todos los dispositivos conectados. Firebase Realtime Database es un claro ejemplo en el mundo móvil: guarda los datos en formato JSON en la nube y los sincroniza automáticamente con cada cliente que esté conectado, sin que tengas que hacer nada complicado de red o polling \cite{FirebaseDocsRTDB}. Esto facilita mucho crear apps colaborativas como chats, dashboards o seguimiento de ubicación en vivo.

% Lo mejor es que funciona bien incluso con internet malo o sin conexión: el SDK de Firebase guarda los datos localmente en el celular, así la app sigue respondiendo aunque estés offline \cite{FirebaseDocsRTDB}. Cuando vuelve la conexión, sincroniza solo lo que cambió de forma automática y resuelve conflictos con reglas que tú defines \cite{FirebaseProductRTDB}. Además, se accede directamente desde el código del móvil y se protege con reglas declarativas, lo que la hace perfecta para apps Flutter donde varios usuarios necesitan ver el mismo estado actualizado en tiempo real.

% --- SECCIÓN 4 ---
\subsection{Integración y Backend}

\subsubsection{Backend as a Service (BaaS)}
Backend as a Service (BaaS) es un modelo en la nube que les da a los desarrolladores de apps móviles un backend ya listo y manejado por otros. Así puedes dejar en manos del proveedor cosas como guardar datos, autenticar usuarios, lógica del servidor y notificaciones push, sin tener que configurar todo desde cero. Esto acelera muchísimo el desarrollo porque te saltas un montón de trabajo manual \cite{Back4appTutorial}. En Flutter hay plataformas como Firebase, Supabase y AWS Amplify que traen SDKs fáciles de integrar, así el equipo se puede enfocar solo en el frontend, la lógica de la app y la experiencia del usuario, sin preocuparse por servidores \cite{TrootechBackendGuide}.

Las ventajas más claras son que reduces un montón el tiempo para lanzar la app (time-to-market), la escalabilidad se maneja sola según crece el uso, y la seguridad suele estar bien cubierta por el proveedor. Entre las opciones más usadas, Firebase es el que más se ve en Flutter por su integración nativa, mientras que Supabase gusta por ser open source y usar PostgreSQL \cite{VoxturrBackend2026}. Es perfecto para prototipos, MVPs o apps que necesitan cosas en tiempo real, aunque para proyectos muy grandes y complejos puede limitar la personalización y te ata un poco al proveedor (vendor lock-in) \cite{TrootechBackendGuide}.

\subsubsection{APIs REST}
Las APIs REST son una forma estándar de diseñar servicios web para que una app (como una hecha en Flutter) pueda hablar con un servidor. Se apoyan en HTTP y usan sus métodos básicos: GET para leer datos, POST para crear, PUT o PATCH para actualizar y DELETE para borrar. Todo se hace sobre recursos que tienen una dirección única (URI) \cite{CodecademyRESTGuide}. En Flutter, para usar una API REST lo normal es hacer peticiones HTTP asíncronas con paquetes como http o dio, y los datos suelen ir y venir en formato JSON \cite{KrasamoBestPractices}.

Para que todo funcione bien y sea fácil de mantener, hay que manejar las operaciones asíncronas con async/await, centralizar las URLs base y los tiempos de espera, y controlar muy bien los errores y los problemas de red \cite{KrasamoBestPractices}. Además, es clave separar la lógica de red en una capa de servicios aparte, usar autenticación con tokens (como Bearer) para las partes protegidas, y crear modelos de datos tipados para que el código sea más seguro y no haya errores raros al procesar las respuestas \cite{MediumAdvancedIntegration}. De esta forma la app queda mucho más ordenada, fácil de probar y de mantener a largo plazo.

\subsubsection{Serialización JSON}
La serialización JSON es básicamente convertir objetos de Dart a texto en formato JSON (para enviar datos) y al revés (para recibirlos), y es algo esencial cuando trabajas con APIs REST o cualquier servicio web en Flutter \cite{FlutterJsonDocs}. Flutter trae la librería dart:convert para hacerlo a mano: tú mismo escribes los métodos toJson() y fromJson() en tus clases modelo. Funciona bien en proyectos chiquitos, pero cuando los modelos se ponen más complejos se vuelve un lío, fácil de equivocarte y difícil de mantener \cite{MediumJsonUnderstanding}.

Para apps medianas o grandes lo mejor es usar serialización automática con generación de código, como el paquete json\_serializable. Le pones anotaciones a tus clases y el paquete genera todo el código de serialización por ti con build\_runner. Así evitas repetir código, reduces errores (porque se detectan al compilar y no cuando la app ya está corriendo), maneja bien tipos complicados, nulos y cambios en los datos, y queda mucho más seguro y fácil de escalar en proyectos grandes \cite{PubJsonSerializable, FlutterJsonDocs}.

% --- SECCIÓN 5 ---
\subsection{Ingeniería de Software y Seguridad}

\subsubsection{Inyección de Dependencias}
La Inyección de Dependencias (DI) es un patrón que evita que una clase cree por su cuenta las cosas que necesita (sus dependencias), y en cambio las recibe desde afuera. Esto hace que el código tenga menos acoplamiento y sea más fácil de mantener. En Flutter es clave para conectar bien las capas de la app (como la vista, el viewmodel, el repositorio o los servicios) sin que todo quede revuelto \cite{FlutterDocComunicacionCapas}. Lo más común es pasar esas dependencias por el constructor de la clase, lo que simplifica un montón las pruebas unitarias (puedes meter mocks fácilmente) y te deja cambiar implementaciones sin tocar el código principal \cite{MediumDIInjector}.

Para manejar todo esto de forma organizada en Flutter se usan contenedores o localizadores de servicios, como get\_it, provider o injectable. Por ejemplo, get\_it es un localizador global y súper ligero donde registras las dependencias: como singletons (una sola instancia para toda la app), fábricas (nueva instancia cada vez que se pide) o lazy singletons (se crea solo cuando se necesita por primera vez) \cite{Q2BInyeccionDependencias}. Así centralizas toda la creación de objetos en un solo lugar, y la lógica de negocio y la UI quedan totalmente separadas de cómo se instancian las cosas.

\subsubsection{Manejo de Errores}
En Dart y Flutter, manejar bien los errores es clave para que la app no se caiga de golpe y el usuario tenga una experiencia decente sin sorpresas feas. Dart usa el clásico try-catch-finally para atrapar excepciones: con on capturas errores específicos que sabes que pueden pasar, y con catch agarras el objeto de la excepción más el StackTrace para ver qué pasó y depurar más fácil \cite{MediumManejoErroresDart}.

Flutter va un paso más allá y tiene sus propios ganchos: FlutterError.onError atrapa errores que pasan al construir o pintar widgets, y para los errores asíncronos que se escapan de ahí se configura PlatformDispatcher.instance.onError \cite{FlutterDocManejoErrores}. Lo ideal es envolver todo en runZonedGuarded para que capture cualquier excepción no manejada en la app, y luego enviar esos reportes a un servicio centralizado como Firebase Crashlytics o Sentry, así sabes rápido qué falla y lo arreglas antes de que afecte a muchos usuarios \cite{MediumArquitecturaManejoErrores}.

\subsubsection{Seguridad de Credenciales}
La seguridad de las credenciales es súper importante porque protege cosas sensibles como contraseñas, tokens de acceso o claves API. Lo primero y más básico es nunca dejar esas credenciales metidas directamente en el código fuente, porque alguien podría hacer ingeniería inversa y sacarlas fácilmente \cite{OWASPMobileTop10}. Si de verdad necesitas alguna en tiempo de compilación, lo mejor es usar variables de entorno o archivos de configuración que nunca subas al repositorio.